\section{Structural Tools for the Millennium Matrix}
\label{app:mm-tools}

\begin{remark}[Purpose of the general MM-tools]
The following tools are completely general. They apply to every admissible
Millennium matrix, independently of whether one later reads that matrix as an
algebraic intertwiner, a closure operator, or a global Heisenberg-compensator
matrix. Their role is structural discipline.
\end{remark}

\begin{definition}[Layer-identity principle]
Let $X$ be a structural object of the Sphaera architecture, and let
\[
\mathcal M_{16}: X \mapsto X^{(\ell)}
\]
be an admissible realization of $X$ on some layer $\ell$. Then $X^{(\ell)}$
is not a new object, but a new realization of the same underlying structural
content.
\end{definition}

\begin{definition}[Projection discipline]
A projection
\[
\Pi: X \mapsto X_{\mathrm{readout}}
\]
is admissible if it only extracts already stabilized structural content. In
particular, projection may reduce visibility, but it may not create new
structure.
\end{definition}

\begin{definition}[Closure-before-readout rule]
Every admissible passage through the Millennium matrix must respect the
ordering
\[
\text{triolic admissibility}
\;\longrightarrow\;
\text{closure compatibility}
\;\longrightarrow\;
\text{projection}
\;\longrightarrow\;
\text{interpretation}.
\]
\end{definition}

\begin{definition}[Sector-stability criterion]
A sector $\Sigma$ is structurally stable under the Millennium matrix if
\[
\mathcal M_{16}(\Sigma)\subseteq\Sigma
\]
or, in the strongest case,
\[
\mathcal M_{16}(\Sigma)=\Sigma.
\]
\end{definition}

\section{Heisenberg-Compensator Tools for Physical Readout}
\label{app:hc-tools}

\begin{remark}[Transition from general MM-tools to HC-tools]
We now impose the stronger reading that the Millennium matrix is not only a
structural intertwiner, but the global matrix realization of the Heisenberg
compensator across admissible layers.
\end{remark}

\begin{definition}[Global Heisenberg-compensator matrix]
The Millennium matrix is called a global Heisenberg-compensator matrix if its
action extends the local compensator rule
\[
C=\tfrac12(\mathrm{id}+J)
\]
to the full admissible multi-layer architecture of the Companion.
\end{definition}

\begin{definition}[Frame-balancing tool]
Let $X=X^+ + X^-$ be decomposed into a frame-even and frame-sensitive part
with respect to an involution $J$, so that
\[
J(X^+) = X^+,
\qquad
J(X^-) = -X^-.
\]
Then the compensator action satisfies
\[
C(X)=\tfrac12(X+JX)=X^+,
\]
and removes the frame-sensitive component.
\end{definition}

\begin{definition}[Fix-sector tool]
A sector $\Sigma_{\mathrm{fix}}$ is an HC-fix sector if every object in it is
invariant under the compensator:
\[
C(X)=X
\qquad\text{for all }X\in\Sigma_{\mathrm{fix}}.
\]
\end{definition}

\begin{definition}[Confinement tool]
The confinement tool is the rule that the Heisenberg-compensator matrix
prevents admissible structural content from escaping the closure-stable sector
selected by the theory.
\end{definition}

\begin{definition}[Physical-readout tool]
A physical regime is an HC-readout if it is obtained as an effective
projection of already compensated and closure-stabilized structural content.
\end{definition}

\section{Notation for the Master Chain}
\label{app:master-chain-notation}

\begin{definition}[Stage symbols of the master chain]
We use the following symbols for the principal stages of the Companion:
\begin{align*}
\mathfrak T &:= \text{triolic source},\\
\mathfrak C &:= \text{carrier layer},\\
\mathfrak C^{\mathrm{adm}} &:= \text{admissibly transported carrier},\\
\mathfrak C^{\mathrm{cl}} &:= \text{closure-stable compensated carrier},\\
\mathfrak G &:= \text{group-layer symmetry readout},\\
\mathfrak L &:= \text{null-axis structure},\\
\mathfrak P &:= \text{photon fix-sector},\\
\mathfrak R &:= \text{residual-localization sector},\\
\mathfrak H &:= \text{Higgs-mediation sector},\\
\mathfrak F &:= \text{field-equation readout}.
\end{align*}
We further write $\mathcal M_{16}$ for the Millennium matrix and
$\mathcal H_{16}$ for its global Heisenberg-compensator reading.
\end{definition}

\begin{definition}[Master chain]
The full master chain is written as
\[
\mathfrak T
\;\longrightarrow\;
\mathfrak C
\;\xrightarrow{\;\mathcal M_{16}\;}
\mathfrak C^{\mathrm{adm}}
\;\xrightarrow{\;\mathcal H_{16}\;}
\mathfrak C^{\mathrm{cl}}
\;\longrightarrow\;
\mathfrak G
\;\longrightarrow\;
\mathfrak L
\;\longrightarrow\;
\mathfrak P
\;\longrightarrow\;
\mathfrak R
\;\longrightarrow\;
\mathfrak H
\;\longrightarrow\;
\mathfrak F.
\]
When intermediate carrier detail is not needed, we also use the short form
\[
\mathfrak T
\to
\mathfrak C
\to
\mathcal M_{16}
\to
\mathcal H_{16}
\to
\mathfrak G
\to
(\mathfrak P,\mathfrak R,J\!\cdot\!\mathfrak R)
\to
\mathfrak H
\to
\mathfrak F.
\]
\end{definition}

\begin{remark}[Why this notation is useful]
This notation is useful because it prevents long later sections from drifting
into disconnected local language. Instead, one may always ask at which stage
of the master chain a given statement belongs.
\end{remark}

\section{Further Technical Appendices}
\label{app:further-technical}



